\section{Introduction}

%% Motivation
The essence of the rank-polymorphic programming model
is implicitly treating all operations as \emph{aggregate} operations,
usable on arrays with arbitrarily many dimensions.
The model was first introduced by Iverson with the language APL \cite{APL}.
Over time, Iverson continued to develop this programming model,
making it gradually more flexible,
eventually leading to the creation of J \cite{J} as a successor to APL.
The boon APL offered programmers
was a notation without loops or recursion:
Programs would automatically follow a control-flow structure
appropriate for the data being consumed.
The nature of the implicit iteration structure
could be modified using second-order operators,
such as folding, scanning, or operating over a moving window.
These second-order operators would directly reveal
all loop-carried data dependences.

In this sense,
other languages demanded that
unnecessary work be put into both compilers and user programs.
The programmer would be expected to
write the program's iteration structure explicitly;
in many languages this entails describing a particular serial encoding
of what is fundamentally parallelizable computation.
The compiler must then perform intricate static analysis
to see past the programmer's overspecified iteration schedule.

The design of APL earned a Turing award for Iverson \cite{IversonTuringLecture}
as well as a mention in an earlier Turing lecture \cite{BackusTuringLecture},
praising it for showing the basis of a solution to the ``ven Neumann bottleneck.''
However APL's subsequent development proceded largely in isolation
from mainstream programming-language research.
The APL family of languages % or just "APL"?
painted itself into a corner with design decisions
such as requiring functions to take only one or two arguments
and making parsing dependent on values assigned at run time.
As a result,
APL compilers were forced to support only a subset of the language
(such as Budd's compiler \cite{BuddCompiler})
or to operate on small sections of code,
alternating between executing each line of the program and compiling the next one
\cite{APL3000Compiler}.
What we gain from
the rank-polymorphic programming model's
natural friendliness to parallelism,
we can easily lose by continually interrupting the program
to return control to a line-at-a-time compiler.
Limiting the compiler to operating over a narrow window of code
can also eliminate opportunities for code transformations like fusion,
forcing unnecessary materialization of large arrays.

The tragedy of rank-polymorphic programming
does not end at forgone opportunities for performance.
Despite the convenience of rank polymorphism
for writing array-processing code%
---a common task in many application domains---%
 % "APL never took off" is historically inaccurate but describes the current state well 
APL and its close descendants do not see widespread use.
There is enough desire for implicitly aggregate computation
to support user communities for systems such as
NumPy \cite{Numpy} and MATLAB \cite{MATLAB},
% I am a bit uncomfortable with making this very nontechnical claim here
which do not follow as principled or as flexible a rule
for matching functions with aggregate arguments%
\footnote{For example, operations which already expect aggregate data%
% This statement combines two objections to prior art:
% 1. In Numpy, user code doesn't get to specify arg rank, so only builtins can
% be rank-polymorphic.
% 2. MATLAB's det function does not lift. If you pass it an argument with 3 or
% more dimensions, it gives an error message demanding a square matrix.
---for example if the programmer writes a function
to compute the norm of a vector or the determinant of a matrix---%
do not always lift easily to consume even higher-dimensional arguments}.
However, programmers are driven away from APL itself by features such as
obtuse syntax,
restrictions on function arity,
poor support for naming things,
and a limited universe of atomic data to populate the arrays \cite{WhatsWrong}.


%% What we're doing
Our goal is to study rank polymorphism itself
without getting bogged down by APL's other baggage.
A formal semantics of rank polymorphism
is the essential groundwork for understanding
how rank-polymorphic programs ought to behave,
how they should be compiled,
and how they can be safely transformed to reduce execution cost.
To that end, we develop Remora,
a language which integrates rank polymorphism with typed $\lambda$-calculus.
% What do we gain from using lambda?

A key problem impeding static compilation of rank-polymorphic programs
is identifying the implicit iteration structure at each function application.
Even without obscuring the programming model
with the idiosyncratic special case behavior APL accreted,
rank polymorphism itself has seemed ``too dynamic'' for good static compilation
due to having its control structure derived from computed data.
The old style of line-at-a-time compilation
relied on inspecting functions and data at run time
to decide what loop structure to emit.

Remora's answer to the problem of finding the iteration space
is a type system which describes the shapes of arrays
and thereby identifies the implicit iteration space for each function application.
In order for types to provide enough detail about array shapes,
we use a restricted form of dependent typing,
in the style of Dependent ML \cite{DependentML}.
In Dependent ML, types are not parameterized over arbitrary program terms
but over a much more restricted language.
For Remora, our language of type indices
consists of natural numbers, describing individual dimensions,
and sequences of natural numbers, describing array shapes.

%% Aside: Our dependent typing is different
Past work on applying dependent types to computing with arrays
has focused on ensuring the safety of accessing individual array elements
\cite{DMLBounds, Qube}.
Bounds checking array indices is essential
in a programming model where extracting a single element
is the only elimination form for arrays,
but the rank-polymorphic programming model generally eschews this operation.
Instead, arrays are consumed whole,
and function application itself serves as the elimination form for arrays.

%% How *we* use dependent types
In Remora, a function's type describes the shapes of the expected argument arrays,
called the ``cells,''
and the type of the atomic data inside the array.
The typing rule for function application
is responsible for identifying the ``frame'' shape,
\ie, the iteration structure derived from the non-cell dimensions.
Type soundness means that
the type system produces more than a safety guarantee:
conclusions it draws about the iteration structure
can be used to correctly compile the program.
Our type system is flexible enough to express polymorphism over the cell shape,
such as a determinant function that can operate on square matrix cells of any size.
It can also handle functions whose output shape is not determined by input shape alone,
such as reading a vector of unknown size from user input
or generating an array of caller-specified shape.


%% What's in this document
We begin with an overview of the rank-polymorphic programming model,
written as a programming tutorial for an untyped variant of Remora.
After developing the intuition for rank polymorphism,
we present a formal description of Remora's core language.
This includes Remora's abstract syntax,
the language of type indices with its associated theory,
the static semantics which identifies array shapes and iteration spaces,
a type-driven dynamic semantics,
and a type-soundness theorem linking the static and dynamic semantics.
Since our formal presentation is intrinsically typed,
we also include an algorithm for partial type erasure,
to characterize which type-level information is truly necessary to keep at run time.
A bisimulation argument connects
the dynamic semantics of explicitly-typed Remora to that of erased Remora.
